\section{3교시 (90분) — 종료: 보고서는 마지막 독자를 위해 쓴다}

% =====================================================================
% 3교시 — 종료 (교재 제4장, 워크북 §3 ⑱)
% =====================================================================
\begin{frame}{3교시에서 배우는 것 — 선언이 아니라 증명}
\begin{itemize}
\item 종료 보고서 8항목 — 모든 숫자가 어느 문서 어느 행에서 왔는지 추적된다
\item BL-02 대비 최종 4축 + 리스크·편익·지속가능성, 원가는 3층을 분리해 적는다
\item 계약 종결 — 검수 \ra{} 정산 \ra{} 유보금 \ra{} 하자보수 \ra{} 지체상금 0의 근거
\item 문서·형상 인수(2023 지적 회수) · 이관 10건 · 잔여 리스크 · 팀 해산
\item 마지막 독자 — 3년 뒤의 내부감사·분쟁·다음 PM
\end{itemize}
\keymsg{종료 보고서는 프로젝트의 마지막 문서가 아니라 내부감사의 첫 문서다}
\note{교재 제4장 도입. 2교시(서비스 전환)에서 SAC 판정 05-21·ELS 종료 04-23이 나왔고, 이 교시는 그 둘을 입력으로 받아 G4(2027-05-28) 종료 보고서를 쓴다. 2019년 세 줄짜리 종료 보고서와 2027년 여덟 항목의 차이는 길이가 아니라 "숫자의 출처가 있는가"임을 먼저 못 박는다. 워크북 §3(⑱)의 항목 1\til{}8이 이 교시의 순서와 같다. 예상 소요 4분.}
\end{frame}

\begin{frame}{종료 보고서의 구조 — 8항목과 입력 목록}
\begin{tbl}[\scriptsize]{L{0.06\textwidth}L{0.40\textwidth}L{0.44\textwidth}}
\textbf{\#} & \textbf{항목} & \textbf{입력(어느 문서에서 오는가)} \\ \midrule
1\til{}2 & 문서 정보·완료 기준 판정 / BL-02 대비 최종 성과 & 헌장 §14 다섯 기준 · 성과보고 제17호(05-31) \\
3\til{}4 & 변경 로그 총괄 / 컨틴전시·관리예비비 집행 총괄 & 변경 로그 v1.4 · 결재·CCB 월 보고 \\
5\til{}6 & 잔여 리스크 이관 / 미해결 사항 이관 목록 & 등록부 v1.4 · SAC 판정 · KEDB · SLA \\
7 & 계약 종결 & 계약 종결 서면 7건 · 검수 확인서 R1\til{}R5 \\
8 & 문서·형상 인수 · 팀 해산 · 승인 & 클린 빌드 05-12 · SBOM · 편익 원장 v1.0 \\
\end{tbl}
\begin{itemize}
\item 첫 문장 — 기간 2026-01-05\til{}2027-05-28(17개월), BL-01 \ra{} BL-02, \textbf{재기준선 0회}
\item 완료 기준 다섯 — 각각 판정·근거·일자·\textbf{판정자 실명}(박준영 SAC · 박세연 형상)
\end{itemize}
\keymsg{추적되지 않는 숫자는 종료 시점에 만들어진 숫자다 — 첫 줄은 입력 목록이다}
\note{교재 4.1. PRINCE2 7 End Project Report(PID 대비 성과·편익 리뷰 계획·교훈·후속 조치)와 PMBOK 8판 Closing Focus Area(성과 평가·계약 종결·재무 정산·기록 보관·팀 해산)를 합친 것이 온담의 8항목이다. 증명해야 하는 것 셋은 2020·2023년 지적사항에서 온다 — 편익 산정 근거, 예비비 집행 승인 근거, 소스·형상 인수. "실무에서 틀리는 것" — 종료 시점에 처음 쓴다(성과보고 17호가 있으면 항목 2는 그 복사다), 완료 기준을 종료 시점에 정한다, 입력 목록이 없다. 워크북 §3.3 항목 1. 예상 소요 8분.}
\end{frame}

\begin{frame}{BL-02 대비 최종 성과 — 네 축, 편차 0·0·$-$5.13·24\%}
\begin{tbl}[\scriptsize]{L{0.08\textwidth}L{0.28\textwidth}L{0.30\textwidth}L{0.12\textwidth}L{0.10\textwidth}}
\textbf{축} & \textbf{BL-02} & \textbf{최종 실적(05-28)} & \textbf{편차} & \textbf{판정} \\ \midrule
범위 & FP 12,585 / RTM 1,189 & 인수 FP 12,585(100\%) & 0 & 이내 \\
일정 & G4 2027-05-28 & G4 05-28 정시, 창 34h·8h 정시 & 0영업일 & 이내 \\
원가 & PMB 146.82 & AC \textbf{151.95} / CPI 0.966 & VAC $-$5.13 & 이내(예외 0) \\
품질 & 0.4건/FP & 1,187건 \ra{} \textbf{0.094건/FP} & 기준의 24\% & 이내 \\
리스크 & 잔여 EMV 8 / 단일 2 & 최대 잔여 7.26(9월) & — & 이내 \\
\end{tbl}
\begin{itemize}
\item 편익 — 첫 측정 B-08 13.6h, 판정은 G5 \ra{} ⑳ 원장으로 이관 · 지속가능성 — 사건 1건(P1-02) 기록
\item 9월 as-is 전망 06-11(+10) \ra{} EX-01 대안 1로 회복 — 0영업일은 회복의 결과다
\end{itemize}
\keymsg{품질을 결함 수로 적지 마라 — 분모 12,585·분자(통합\til{}UAT)·기준 0.4가 있어야 1,187이 읽힌다}
\note{교재 4.2 표. 여섯 열(기준선/실적/편차/허용범위/판정/참고 BL-01) 중 슬라이드는 다섯을 보인다. 범위 +185 FP는 BL-01 대비이며 허용 ±372의 49.7% — 재확인 트리거 ±3% 미발동. 품질 총 발견 2,563(설계 210·스프린트 1,166·통합\til{}UAT 1,187)이고 밀도의 분자는 통합시험\til{}UAT만이다 — 9월 전망 0.09 적중. 리스크 여유 최소 0.74(11월). 워크북 §3.3 항목 2. 예상 소요 8분.}
\end{frame}

\begin{frame}{원가 3층 — "16.05 미사용"과 "VAC $-$5.13"이 동시에 참}
\fig[0.60]{../그림/PM_Day4/fig_4_1_three_layers_final.pdf}
\figcap{그림 4-1. 승인 168.00 / 집행 한도 156.00 / PMB 146.82 대 최종 AC 151.95. EAC 이력 9호 150.86 \ra{} 11호 152.40(경고) \ra{} 17호 151.95, 예외 진입 0회}
\keymsg{PM의 자는 PMB다 — 종료 보고서는 세 층을 분리해 적고 "절감"이라는 말을 쓰지 않는다}
\note{교재 4.2 그림 4-1. 승인 예산에는 스폰서만 쓰는 관리예비비 12.0이 들어 있어 처음부터 PM의 돈이 아니고, 한도와 PMB 사이의 컨틴전시는 리스크가 실현되면 쓰라고 둔 돈이다. 세 줄 — 승인 대비 16.05 미사용(관리예비비 12.0 + 컨틴전시 반납 4.05), 한도 대비 4.05 여유, PMB 대비 VAC $-$5.13. 강조할 것: 9월 조정 CPI 0.968이 최종 0.966을 0.002 안에서 예측했다 — Day3에서 착시 1(라이선스 4.70)을 걷어 냈기 때문이며, 걷어 내지 않은 CPI 1.032로 EAC1 141.23을 채택했다면 최종은 "예상 못 한 10.7억 초과"로 기록됐을 것이다. 워크북 §3.3 항목 2·4. 예상 소요 8분.}
\end{frame}

\begin{frame}{변경·예비비 총괄 — CR-01\til{}14와 컨틴전시 10.20의 분해}
\begin{columns}[T]
\begin{column}{0.46\textwidth}
\subhead{변경 로그 총괄}
\begin{itemize}
\item 승인 10 · 조건부 1(CR-09) · 거버넌스 1(CR-10) · 기각 1(CR-05) · 이관 1(CR-12)
\item Day3 이후 — CR-12 P1 범위 외 이관 · CR-13 R-10 재정제 0.45 · CR-14 패치 0.20
\item FP 누적 +185(+1.5\%, 허용의 49.7\%), SI +1.02
\end{itemize}
\end{column}
\begin{column}{0.52\textwidth}
\subhead{컨틴전시 10.20 = 다섯 조각}
\begin{itemize}
\item 편입 1.02(변경 대가 \ra{} PMB) + 리스크 배정 집행 2.50
\item 미배정 풀 0.77 + 편차 흡수 1.86 + \textbf{반납 4.05}
\item 소진율(편입 제외) 50.3\% · 관리예비비 12.0 요청 0·\textbf{반납 12.0}
\item 유보 3건 — 클라우드 3.0 반납 / PMO 3.0 승인 전 / 전환 2.0 감액
\end{itemize}
\end{column}
\end{columns}
\keymsg{모든 집행에 결재(PM + 재경)와 CCB 월 보고가 붙어 있다 — 윤태호의 질의에 이 표가 답이다}
\note{교재 4.2 후반, 그림 4-2(보조). 리스크 배정 집행 2.50의 내역 — R-07 0.15·R-05 0.42·R-03 0.50·R-10 0.45·R-19 0.60·R-20 0.18·R-13 0.20. 미배정 풀 0.77 — 단말 0.30·아카이브 0.20·부하 시나리오 0.20·보안 0.05·임시 라이선스 0.02. 편차 흡수 1.86은 리스크 배정 없는 AC 초과(클라우드 0.76·망분리·DR 0.50·PMO 0.60)이며 이것이 VAC의 본체다. 관리예비비 대상 R-01·R-14·R-17 미실현(R-01은 분리 의결로 흡수). "컨틴전시를 합계로만 적는다"가 가장 흔한 오류 — 어느 리스크에 어느 결재로 나갔는지가 없으면 내부감사 질의에 답할 수 없다. 워크북 §3.3 항목 3·4. 예상 소요 8분.}
\end{frame}

\begin{frame}{계약 종결 — 검수 \ra{} 정산 \ra{} 유보금 \ra{} 하자보수 \ra{} 지체상금}
\begin{tbl}[\scriptsize]{L{0.14\textwidth}L{0.50\textwidth}L{0.28\textwidth}}
\textbf{단계} & \textbf{SI 한빛디지털} & \textbf{근거 서면} \\ \midrule
검수 & 최종 검수 확인서 05-26 & 검수 확인서 R1\til{}R5 \\
정산 & 69.22 + 대기 0.60 + 컨틴전시 귀속 1.20 = \textbf{71.02}, 잔금 05-31 & 합의서 7호(05-26) \\
유보금 & 5\% = 3.46억 \ra{} 하자보수보증증권 대체 & 계약 §16 · 보증증권 \\
하자보수 & 1년 2027-05-29\til{}2028-05-28, 범위 = \textbf{인수 산출물의 결함} & 개선·요구 변경은 P5 \\
지체상금 & \textbf{0 — 근거 네 줄}(다음 프레임) & 합의서 6호 · Day3 ⑮ 항목 6 \\
\end{tbl}
\begin{itemize}
\item 나머지 — 전환 용역 10.05(하자 6개월) · 상용SW 19.42 · 클라우드 22.16 중 P5 승계 · 단말 8.90
\item 하도급(모빌커넥트 220 FP) — 수행사 정산 내 1.21, \textbf{직접 지급 요청 없음} 확인
\end{itemize}
\keymsg{정산액을 한 숫자로 적지 마라 — 71.02가 69.22 + 0.60 + 1.20임을 모르면 0.60은 "왜 줬나"가 된다}
\note{교재 4.3. 순서가 곧 논리다 — 검수 없이 정산 없고, 정산 없이 유보금 없다. 하자보수 범위 정의가 없으면 하자보수 1년은 무상 유지보수 1년이 된다(결함 vs 개선의 판정은 CAB). 하도급 직접 지급 여부를 확인하지 않으면 종료 후 정산이 다시 열린다. 상용SW 잔금 4.70은 통합시험 통과 01-15에 발생 — Day3 착시 1의 해소 시점. 워크북 §3.3 항목 7. 예상 소요 8분.}
\end{frame}

\begin{frame}{지체상금 0의 네 줄 — 같은 표가 대기 비용 0.60도 설명한다}
\begin{itemize}
\item ① G4 05-28 = 기준선 완료일 \ra{} 계약 §18의 "지연" 불성립
\item ② 순 지연 분해(Day3 ⑮ 항목 6) — 발주 7 + 동시 3 + 수행 3 $-$ 회복 3 = 10 \ra{} EX-01로 0
\item ③ 수행사 단독 3일(S10 재작업) \ra{} 회복 3일(S12\til{}S17)과 상계 — 변경협의체 10-14
\item ④ 동시 3일 — 07-15 합의(SCL 원칙 10) 공기 연장 인정·비용 불인정
\item 같은 귀책 분해표의 반대 방향 = 발주사 귀책 대기 비용 0.60 지급(R-05 0.42·R-20 0.18)
\end{itemize}
\keymsg{네 줄이 없으면 0은 "봐줬다"로 읽히고, 있으면 "계산했다"로 읽힌다 — Day3 동시기록이 종료에서 정산된다}
\note{교재 4.3 그림 4-3(보조). 0은 "없음"으로 읽히기 쉬운 숫자이므로 근거가 필요하다 — 내부감사는 "왜 물리지 않았는가"를 묻는다. 핵심 문장 — 같은 표가 양쪽(지체상금 0과 대기 비용 지급)을 다 설명할 때만 그 표는 정산의 근거가 된다. 이것이 Day3 5.5절 동시기록 문화의 종료 시점 회수다. 귀책을 그때그때 적지 않았다면 종료 시점에 이 네 줄은 쓸 수 없다. 워크북 §3.3 항목 7. 예상 소요 6분.}
\end{frame}

\begin{frame}{문서·형상 인수 — 2023년 지적사항의 회수}
\begin{itemize}
\item 조건은 "받았다"가 아니라 \textbf{발주사 저장소에서 클린 빌드가 재현되는가}
\item 온담 — 재현 100\%(05-12, 12개 모듈 + 하도급 어댑터 220 FP), SBOM 312(GPL 계열 0)
\item 빌드·배포 절차서 v1.0 · 일 1회 동기화 \ra{} G4 후 발주사 저장소 단독 정본
\item 아카이브 \textsf{governance/p1-closure/} — 헌장\til{}편익 원장, 계약 서면 32건, 보존 10년
\item 지적 대응 요약 — 2020 편익 근거 \ra{} ⑳ 원장 / 2020 손실 확정 \ra{} §14 / 2023 소스 \ra{} 형상 인수
\end{itemize}
\keymsg{종료 보고서는 지적사항에 답하는 자리이지, 지적사항이 없었던 척하는 자리가 아니다}
\note{교재 4.4. 소스 파일을 받는 것은 2023년 지적("위탁 산출물 형상관리·소스 인수 절차 미비")의 상태 그대로다. 형상 인수 = SAC 항목 9의 증빙이자 헌장 §14 ②의 판정 근거. SBOM이 없으면 라이선스 리스크가 운영 중에 발견된다. 실무 오류 — 하도급분이 빠진다, 아카이브 위치가 PM의 노트북이다. 워크북 §3.3 항목 8. 예상 소요 6분.}
\end{frame}

\begin{frame}{미해결 사항 이관 목록 — 소유자 실명 없는 이관은 소멸이다}
\begin{tbl}[\scriptsize]{L{0.04\textwidth}L{0.42\textwidth}L{0.18\textwidth}L{0.14\textwidth}L{0.12\textwidth}}
\textbf{\#} & \textbf{내용} & \textbf{소유자} & \textbf{기한} & \textbf{승인} \\ \midrule
1 & 아카이브 접근 절차 재경 승인(SAC 11 조건부) & 재경팀장 & 06-15 & 승인 \\
3 & 발주 마감 22:00 운용 전환, CAB 재상정(CR-09 ①) & 한동석 & 2027-09 CAB & \textbf{승인 전} \\
6 & 이천·안성 검사기록 디지털화 — 편익 B-02 분모의 전제 & 나영선 & 09 이사회 & \textbf{승인 전} \\
7 & SAP ECC 2027-12 EoM 대응, IF-SAP 8본(KE-04) & 박세연 & 08-31 & \textbf{승인 전} \\
10 & 하자보수기 PMO 잔류 1명 예산 1.05(Day2 유보 3.0) & 프로그램 PMO장 & 06-01 운영위 & \textbf{승인 전} \\
\end{tbl}
\begin{itemize}
\item 10건 = 승인 6 · 승인 전 4 — 세 칸이 결정적: \textbf{소유자 실명 · 기한 · 승인 여부}
\item 출처 열이 앞선 열아홉 문서의 어느 행이 프로젝트 밖으로 나가는지를 기록한다
\end{itemize}
\keymsg{"승인 전"은 종료 보고서에서 가장 짧은 줄이지만 가장 오래 남는 줄이다}
\note{교재 4.5 표(10건 중 5건 발췌 — 승인 전 4건 전부 + 1번). 나머지 승인 5건: 2 미동의 9점 2차 배포(오정근 07-31) · 4 3PL 재협상(한동석 06-30) · 5 앱 벤더 계약 갱신(오정근·박세연) · 8 클라우드 사이징·유보 3.0 반납(박준영·재경팀장) · 9 ISMS-P 증적(최영주). 이관 항목은 두 종류다 — 소유자와 예산이 있어 실행만 남은 것(승인)과, 소유자는 있으나 결정이 없는 것(승인 전). 후자를 "승인"처럼 적으면 이사회 통과 순간 아무도 책임지지 않는 항목이 된다. 6번을 특히 짚는다 — 범위 밖이지만 편익의 전제이므로 적는다. 적지 않으면 G5 편익 미달이 "P1 결함"으로 기록된다. 워크북 §3.3 항목 6. 예상 소요 8분.}
\end{frame}

\begin{frame}{이관 10건의 출처와 수신 — 다른 넷과 묶는 열}
\fig[0.60]{../그림/PM_Day4/fig_4_4_followon_map.pdf}
\figcap{그림 4-4. 출처(SAC·SLA·KEDB·변경 로그·컨틴전시 유보·편익 원장·정본 제외 범위) \ra{} 항목·소유자·기한 \ra{} 수신 조직. 승인 6 실선·승인 전 4 파선, 6번 \ra{} ⑳ BR-05}
\keymsg{출처 열이 없으면 종료 시점에 누군가 떠올린 목록이고, 있으면 앞선 문서의 기록이다}
\note{교재 4.5 그림 4-4. 어디에서 오는가 — SAC 조건부(1), SLA 강등 재상정 조건(4), KEDB 의존·협의(2·7), 변경 로그 이관·조건부(3), 컨틴전시 유보(8·10), 편익 원장 전제(6), 회사 정본 제외 범위(7). 수신 열에는 결정 회의체(CAB·프로그램 이사회·운영위)가 적혀 있다 — 종료 보고서가 이사회를 통과하는 것과 승인 전 4줄이 결정되는 것은 별개다. 실무 오류 네 가지 — 소유자가 부서명, 승인 여부 열 없음, 범위 밖은 적지 않음, 출처 열 없음. 예상 소요 5분.}
\end{frame}

\begin{frame}{잔여 리스크의 운영 이관과 팀 해산 — 사람의 전환}
\begin{columns}[T]
\begin{column}{0.50\textwidth}
\subhead{잔여 리스크 — 등록부 v1.4(G4 종결)}
\begin{itemize}
\item 종결 12 / 운영 이관 3 / 편익 이관 1
\item R-13 \ra{} KE-01(오정근) · R-16 \ra{} KE-02(한동석) · R-18 \ra{} 앱 벤더(오정근·박세연)
\item R-12 \ra{} BR-01(재무본부·한동석·나영선)
\item 신규 운영 리스크 3 — 3PL 폴백 영구화·SAP EoM·IPO 변경 동결 창
\end{itemize}
\end{column}
\begin{column}{0.48\textwidth}
\subhead{팀 해산 — 사람별 일자}
\begin{itemize}
\item 발주사 14명 — 12명 06-01 복귀, \textbf{2명 P5 전환}
\item 수행사 05-31 철수, 하자보수 대기 4명 비상주
\item PMO 3명 — 2명 프로그램 PMO, 1명 잔류(승인 전, 이관 10)
\item 하자보수 창구 — P5 서비스운영팀(박준영) \ra{} UC ②, 판정 CAB
\end{itemize}
\end{column}
\end{columns}
\keymsg{리스크는 사라지지 않는다 — 운영 등록부·KEDB·편익 리스크 표 중 한 곳에 새 소유자와 함께 다시 적힌다}
\note{교재 4.6. 세 곳 어디에도 적히지 않은 잔여 리스크는 소멸이고, 실현될 때 "아무도 몰랐던 문제"로 돌아온다. 신규 운영 리스크 3건은 운영 리스크 등록부(박준영)로, 검토는 월 1회 서비스 리뷰. 2명의 P5 전환은 SAC 항목 5(교육)의 다른 얼굴 — 만든 사람이 운영 조직에 남는 것이 가장 짧은 지식 이전이다. 실무 오류 — 잔여 리스크 일괄 "종결", 해산을 한 날짜로, 하자보수 창구가 없어 첫 하자 문의가 전 PM의 개인 전화로, 잔류 인력 예산 미결정 채 잔류. Bridges의 "끝냄"은 팀에도 필요하다. 워크북 §3.3 항목 5·8. 예상 소요 7분.}
\end{frame}

\begin{frame}{문서의 마지막 독자 — 3년 뒤의 내부감사·분쟁·다음 PM}
\begin{tbl}[\scriptsize]{L{0.16\textwidth}L{0.38\textwidth}L{0.38\textwidth}}
\textbf{독자} & \textbf{묻는 것} & \textbf{답하는 곳} \\ \midrule
내부감사(3년 뒤) & 컨틴전시 R-03 0.50은 누가 언제 승인했나 & 항목 4 결재 08-12 · CCB 10-15 소급 · 11월\til{} 자동 안건 \\
분쟁(하자보수기) & 이것이 결함인가 개선인가 · 지체상금 0은 봐준 것인가 & 하자보수 범위 정의 · 검수 확인서 R1\til{}R5 · 네 줄 \\
다음 PM(P2·P5) & 온담은 컷오버를 어떻게 했는가 & 런북 실행 기록 열 · 교훈 L-15 \\
\end{tbl}
\begin{itemize}
\item 약어를 풀어 쓴다 — RB-3이 무엇인지 3년 뒤에는 아무도 모른다
\item 출처를 적는다 — 어느 문서 어느 행
\item \textbf{판단의 근거를 결론과 분리해 적는다} — "충족"이 아니라 "조건부 충족, 06-15, 재경팀장"
\end{itemize}
\keymsg{결론만 남은 문서는 3년 뒤 검증할 수 없고, 검증할 수 없는 문서는 신뢰되지 않는다}
\note{교재 4.7\til{}4.8. 06월 프로그램 이사회는 첫 독자이지 마지막 독자가 아니다. 이 장의 오류는 전부 "선언과 증명을 구별하지 못하는 것"으로 수렴한다 — 종료 시점에 처음 쓰고, 원가를 한 층으로("16억 절감"), 지체상금 0을 "없음"으로, 소스를 "받았다"로, 소유자를 부서명으로, 잔여 리스크를 일괄 종결, 해산을 날짜 하나로. 그 결과는 2019년 세 줄과 같은 종류의 문서이고, 마지막 독자는 그것을 처음부터 다시 만들어야 한다. 예상 소요 6분.}
\end{frame}

\begin{frame}{3교시 핵심 정리 — 종료 보고서는 증명이다}
\begin{itemize}
\item 8항목 · 모든 숫자에 출처 · 완료 기준 다섯에 판정자 실명 · 첫 문장 "재기준선 0회"
\item 원가 3층 분리 — 16.05 미사용 / 여유 4.05 / VAC $-$5.13 동시에 참, PM의 자는 PMB
\item 정산 71.02 = 69.22 + 0.60 + 1.20 · 지체상금 0의 네 줄 · 클린 빌드 100\%·SBOM 312
\item 이관 10건(승인 6·승인 전 4) — 소유자 실명·기한·승인 여부 없으면 소멸
\item 잔여 리스크는 새 소유자와 함께 다시 적힌다 · 해산은 사람별 일자
\end{itemize}
\keymsg{4교시 실습 ① — 워크북 §3 ⑱ 종료 보고서·이관 목록·계약 종결을 온담 값으로 직접 쓴다}
\note{제4장 핵심 정리 다섯 줄. 4교시 실습으로 넘기기 전에 워크북 §3.4 빈 템플릿과 §3.5 완성 예시본의 값이 이 교시의 숫자와 같음을 확인시킨다(교재와 워크북이 다르면 워크북이 정본). 질문 시간 포함. 예상 소요 8분.}
\end{frame}

